% Drafted with an AI assistant (limitations and deviations on 31 Oct, implications on 10 Nov 2026).
% Revise everything in your own words. Directional words are tagged "% claim:<id>" and are
% re-checked by `uv run qdrift paper`.
\section{Discussion}
\label{sec:discussion}

\textbf{Implications.} A calibration emulator is a useful screen but not a stand-in for
hardware. It ordered our configurations well yet overstated their accuracy by \val{h1-mean-pp}
points on average. % claim:emu-optimistic claim:emu-ranks-well
A study of calibration-aware QAOA compilation on Heron processors also found that emulators
overstated the benefit measured on hardware~\cite{qaoa2026drift}, so the bias is unlikely to be
specific to classifiers. Day-to-day drift was real but differed between devices. % claim:drift-one-device
A result from one session, such as the single day of hardware data in the adaptive conformal
study~\cite{spencer2026aqcp}, can therefore be off by several points, while repeating a run on
another day costs about \val{session-usage-mean}~s here.

Under a free budget the cheap repairs buy little. The head trained on emulated noisy features
(R4) is free and gave the largest gain and the best calibration, % claim:r4-best claim:r4-best-ece
and measured readout mitigation is the only paid repair that helped. Normalisation, part of the
QuantumNAT pipeline that raised accuracy by up to 43\% on IBM devices of
2021~\cite{wang2022quantumnat}, gave nothing here. % claim:r1-no-gain
On these processors the global scale and shift of the features is modest and stable, so the
remaining loss comes from errors that differ between inputs and qubit chains and that one global
correction cannot remove; the association of feature error with the CZ error and $T_2$ of each
chain points the same way. % claim:h4-not-readout claim:t2-negative
For uncertainty, the practical fix is to calibrate on hardware: labels from any session
restored coverage, and ACI kept it at the target with smaller sets. % claim:hwcal-over claim:aci-near

\textbf{Limitations.} On the free plan a job runs when the queue allows, not at a chosen time:
queue waits reached \val{queue-max-h}~h during the campaign (\val{queue-max-all-h}~h including
the pilots), and with one job per session the campaign holds \val{sessions-done} sessions. The
models are small (\val{n-qubits} qubits, at most \val{max-cz} CZ gates) and the study covers two
datasets with \val{n-test} test inputs per session, so the accuracy of a single session is
uncertain by several points and only effects pooled over sessions are resolved. The qubit chains
were chosen from one day's calibration and then kept fixed; this pairs the data across sessions
but means that later sessions may not use the best qubits of that day. The emulator treats each
chain on its own and contains no crosstalk, and the packing check used \val{pilot-samples} inputs
per device, enough to exclude only large crosstalk effects. Because IBM refreshes the reported
calibration often, its age at job start lay between \val{cal-age-min-h} and \val{cal-age-max-h}~h,
which makes calibration age a weak measure of drift. Finally, twirling and dynamical decoupling
were switched off by design, so the results describe raw device noise; we did not test the
repairs together with these techniques.

\textbf{Deviations from the registration.} The registration is the git tag \texttt{prereg-v1}
(\val{prereg-date}). Before the hypotheses were first evaluated on hardware data, we corrected
the analysis code so that calibration age is measured at job start instead of at the time the
snapshot was downloaded; only the age covariate of H4 is affected. The registration asked for
sessions at a similar time of day, but the queue and the author's schedule did not always allow
it: \val{sessions-day} sessions started in the daytime and \val{sessions-night} at night (India
time). \val{window-sentence} \val{sessions-sentence}
